\DocumentMetadata{
  pdfversion=2.0,
  pdfstandard=ua-2,
  lang=en-US,
  tagging=on,
  tagging-setup={math/setup=mathml-SE,tabsorder=structure}
}
\documentclass[11pt,letterpaper]{article}
\usepackage[margin=0.68in,includefoot]{geometry}
\usepackage{newcomputermodern}
\usepackage{amsmath,amscd,mathtools}
\usepackage{tikz,adjustbox}
\providecommand{\square}{\mdlgwhtsquare}
\usepackage[hidelinks]{hyperref}
\hypersetup{
  pdftitle={Algebra First-Year Exam, August 2016, Part 1},
  pdfauthor={University of Florida Department of Mathematics},
  pdfsubject={Graduate mathematics examination},
  pdfdisplaydoctitle=true,
  bookmarksopen=true
}
\setcounter{secnumdepth}{0}
\setlength{\parindent}{0pt}
\setlength{\parskip}{0.28em}
\setlength{\emergencystretch}{3em}
\setlength{\leftmargini}{2.15em}
\setlength{\leftmarginii}{2.65em}
\setlength{\leftmarginiii}{3em}
\pagestyle{empty}
\raggedbottom
\allowdisplaybreaks
\NewTaggingSocketPlug{block/list/label}{examproblem}{%
  \tagstructbegin{tag=Lbl}\tagstructend%
  \tagstructbegin{tag=\LBody}%
  \tagstructbegin{tag=\ProblemHeadingTag,title-o={\ProblemHeadingText}}%
  \tagmcbegin{tag=\ProblemHeadingTag,actualtext={\ProblemHeadingText}}%
  #2%
  \tagmcend%
  \tagstructend%
}
\newenvironment{examproblems}[2]{%
  \def\ProblemHeadingTag{#2}%
  \AssignTaggingSocketPlug{block/list/label}{examproblem}%
  \begin{enumerate}%
  \setcounter{enumi}{#1}%
  \renewcommand{\labelenumi}{\textbf{\theenumi.}}%
  \setlength{\topsep}{0.35em}%
  \setlength{\partopsep}{0pt}%
  \setlength{\itemsep}{0.48em}%
  \setlength{\parsep}{0.18em}%
}{\end{enumerate}}
\newenvironment{examsubparts}{%
  \AssignTaggingSocketPlug{block/list/label}{default}%
  \begin{enumerate}%
  \setlength{\topsep}{0.18em}%
  \setlength{\partopsep}{0pt}%
  \setlength{\itemsep}{0.16em}%
  \setlength{\parsep}{0.1em}%
}{\end{enumerate}}
\newenvironment{examinstructions}{%
  \par\begingroup\itshape
}{%
  \par\endgroup\vspace{0.18em}
}
\makeatletter
\renewcommand\subsection{\@startsection{subsection}{2}{0pt}%
  {-1.25ex plus -.25ex minus -.1ex}%
  {0.55ex plus .1ex}%
  {\normalfont\large\bfseries}}
\renewcommand{\@maketitle}{%
  \newpage\null\vskip -1em%
  {\footnotesize\bfseries
    \href{https://www.ufl.edu/}{University of Florida}, \href{https://math.ufl.edu/}{Department of Mathematics}\par}%
  \vskip -0.5em%
  \tagstructbegin{tag=H1,title={Algebra First-Year Exam, August 2016, Part 1}}%
  \tagpdfparaOff%
  \tagmcbegin{tag=H1}%
    {\LARGE\bfseries\href{https://gma.math.ufl.edu/exams/algebra-first-year/}{Algebra First-Year Exam}, August 2016, Part 1\par}%
  \tagmcend%
  \tagpdfparaOn%
  \tagstructend%
  \vskip 0.25em%
  \tagmcbegin{artifact}%
    \hrule height 1pt%
  \tagmcend%
  \vskip 1em%
}
\makeatother
\title{Algebra First-Year Exam, August 2016, Part 1}
\author{}
\date{}
\begin{document}
\maketitle
\thispagestyle{empty}
\begin{examinstructions}
Answer four problems. (If you turn in more, the first four will be graded.) Within reason, you may quote theorems as long as you state them clearly.
\end{examinstructions}
\begin{examproblems}{0}{H2}
\def\ProblemHeadingText{Problem 1.}
\item
Let~\(G\) be a cyclic group, and let~\(H\) be a subgroup of~\(G\). Prove that~\(H\) is cyclic.
\def\ProblemHeadingText{Problem 2.}
\item
Let~\(G\) be a finite group, and suppose it is acting on the finite set~\(\Omega\). Suppose~\(N\) is a normal subgroup of~\(G\) and~\(N\) acts transitively on~\(\Omega\). Let \(\omega \in \Omega\), and let \(G_\omega\) be the stabilizer in~\(G\) of~\(\omega\). Prove that \(|G:N|=|G_\omega:G_\omega\cap N|\).
\def\ProblemHeadingText{Problem 3.}
\item
Let~\(G\) be a finite group, let~\(H\) be a subgroup of~\(G\), and let \(n=|G:H|\) with \(n>1\).
\begin{examsubparts}
\item[\textnormal{(a)}]
Prove that if~\(G\) is simple then~\(G\) is isomorphic to a subgroup of the symmetric group \(S_n\).
\item[\textnormal{(b)}]
Give an example to show that, under the general assumptions of the question, if~\(G\) is not simple it is possible that~\(G\) is not isomorphic to any subgroup of \(S_n\) (and prove that your example has these properties).
\end{examsubparts}
\def\ProblemHeadingText{Problem 4.}
\item
Let~\(p\) be a prime, and let~\(G\) be a finite group such that each element of~\(G\) has order a power of~\(p\). Prove that, if~\(G\) is not trivial, then the center \(Z(G)\) of~\(G\) is not trivial. Deduce that~\(G\) is nilpotent.
\def\ProblemHeadingText{Problem 5.}
\item
Let \(S_n\) be the symmetric group in \(n\geq 1\) letters. Let \(\sigma\in S_n\). Describe what is meant by the cycle type of~\(\sigma\). Prove that two elements \(\sigma,\tau\in S_n\) are conjugate to each other in \(S_n\) if and only if they have the same cycle type.
\def\ProblemHeadingText{Problem 6.}
\item
Prove that a group of order \(30\) must have a normal subgroup of order \(15\).
\end{examproblems}
\end{document}
